\documentclass{sample}
\begin{document}
More on displayed formulas with LaTeX. Authors can typset displayed formulas with \textbackslash \{equation\} (known as an equation environment), and furthermore number with \textbackslash begin\{equation\} and internally label (reference) equations with \textbackslash label\{E:uniqueReference\}, as shown below:
\begin{equation}\label{E:integralOne}
\int_0^{\pi} \sin x \, dx = 2
\end{equation}

The reference in the typsetted document can be demonstrated with \textbackslash ref and \textbackslash pageref: 

``Please refer to~(\ref{E:integralOne}) on page~\pageref{E:integralOne}.''

Note that the source file will need to be typset \textbf{twice} in order to get the references updated. The use of \texttt{E:} is only a convention for \textbf{E}quations.

It is possible to build equations without numbering by introducing an asterisk, as \textbackslash begin\{equation\textasteriskcentered\}. This would be equivalent to \textbackslash [. Indeed, it may prove more useful to do this for documents with a large number of equations, where the author then decides which to number and which not to, by deleting the asterisk.

The above equation numbers are updated automatically by LaTeX, as \textit{relative} numbering. It is also possible to assign an \textit{absolute} numbering, as a tag, with \textbackslash tag\{absoluteId\}:
\begin{equation}
\int_0^{\pi} \sin x \, dx = 2 \tag{intOne}
\end{equation}

The alignment of equations for derivations can be realised with \textbackslash begin\{align\} (known as the align environment). The \textit{alignment point} is applied with \& and the line separator is given by two backslashes.
\begin{align*} 
2x - 5y &=  8 \\ 
3x + 9y &=  -12
\end{align*}

Removing the asterisks will yield numbered equations:
\begin{align} 
2x - 5y &=  8 \\ 
3x + 9y &=  -12
\end{align}

Finally, leave out numbering for select equations with \textbackslash notag:
\begin{align} 
2x - 5y &=  8 \notag \\ 
3x + 9y &=  -12
\end{align}

To annotate a line with two ampersands (e.g. for working):
\begin{align}
x &= x \wedge (y \vee z) &&\text{(by distributivity)}
\end{align}

Cases can be typset like a matrix with \textbackslash begin\{cases\} inside an equation or align environment as follows:
\[
f(x)=
\begin{cases}
-x^2, &\text{if $x < 0$;}\\
\alpha + x, &\text{if $0 \leq x \leq 1$;}\\
x^2, &\text{otherwise.}
\end{cases}
\]
\end{document}